\section{Samma frågor, två svar}
\label{sec:sameq}

Hela boken kretsar kring fem frågor. Första delen besvarade dem i mjukvara, andra delen lät
hårdvaran göra det:

\begin{booktable}{@{}lLL@{}}
\thead{Fråga} & \thead{Eget protokoll} & \thead{CAN} \\ \midrule
Var börjar meddelandet? & SOF, två bytes & SOF-bit; sju recessiva bitar som slutmärke \\
Hur långt är det? & LEN, en byte & DLC, fyra bitar, 0--8 databytes \\
Vem är det till? & DST och SRC & \textbf{Ingen adress.} Identifieraren beskriver innehållet \\
Kom det fram helt? & Summacheck\-summa & CRC-15, plus bitstoppning och formatregler \\
Kom det fram alls? & ACK, timeout, omsändning, dubblett\-detektering & En ACK-bit inuti framen \\
\end{booktable}

Att kunna fylla i den tabellen ur minnet, och motivera varje rad, är ungefär vad duggan kräver av
den första halvan.

\section{Vad duggan täcker}
\label{sec:exam}

\subsection{Del A --- det egna protokollet}

Framedesign och checksumberäkning för hand. Parsning som tillståndsmaskin. Adressering, routing och
sekvensnummer. Byte-strömmar som ska tolkas och analyseras. Felmodellen: tappade bytes, korrupt
data, fördröjning. Kedjan ACK $\rightarrow$ timeout $\rightarrow$ omsändning $\rightarrow$
dubbletter.

\subsection{Del B --- CAN och drivrutinen}

CAN-ID, prioritet, DLC och arbitrering. Vilka mekanismer hårdvaran utför, och var. Registerkartan:
vad varje register betyder, och hur varje statusbit nollställs. Databyteordningen över
\code{TX_DATA_LO}/\code{HI}. SPI-transaktionen: kommandobyten, de fyra databyten, \code{SS}-reglerna,
och att koda och avkoda en transaktion för hand. Lagerarkitekturen. De tre begränsningarna.

Registerkartan delas ut med duggan --- indexen behöver inte memoreras. Att \emph{läsa} kartan och
veta vad den betyder behövs desto mer.
